% GUID v6, Part III: the precision Standard Model statements SM1 to SM5 against data. Numbers: docs/simulation/standard_model.py,
% results in docs/simulation/results/standard_model.{json,md}; \chk{sm:SMx.y} names a row of that file, \chk{data:Dn.m} a row of
% predictions_vs_data.json. Statements and scoring: docs/simulation/sm_preregistration.md (fixed before the script existed).
% Home of: alpha(m_Z), the Higgs-coupling slope, 0vbb and the Dirac step, R = dln mu / dln alpha, the coupling crossings.
\section{The Standard Model statements against precision data}\label{sec:sm}

This section tests Part~I's five statements about the particles and couplings of the Standard Model, SM1 to SM5
(Sections~\ref{sec:fields}, \ref{sec:cosmo}, \ref{sec:alpha} and~\ref{sec:matter}), as a preregistration fixed them before any analysis;
Primer~\ref{box:particle} gives the particle physics they use.

\subsection{Statements and data}\label{sec:sm-method}

The preregistration (\texttt{docs/simulation/sm\_preregistration.md}) states the five statements and fixes, for each, the data and the verdict.
\begin{itemize}[nosep]
\item SM1, the spin ladder: elementary quanta have spin 0, 1/2 or 1 (Section~\ref{sec:matter}).
\item SM2, Hypothesis~1$'$ read for matter: the Higgs vacuum grants every elementary rest mass (the \emph{grant}). Higgs couplings are then
proportional to mass, neutrinos are Dirac particles, the photon and gluons are massless, and $m_W=gv/2$ (Section~\ref{sec:fields}).
\item SM3, Hypothesis~\ref{h:info} with the action~\eqref{eq:action}: the information scalar multiplies only the Maxwell term
(Section~\ref{sec:alpha}).
\item SM4, Hypothesis~6$'$: within causality the forces were never one (Section~\ref{sec:cosmo}).
\item SM5, the electroweak split: $\alpha$ is born near $1/128$ at the Z mass and runs to $1/137$ (Section~\ref{sec:matter}).
\end{itemize}
We score every statement as written. Where one turned out to be ill-posed, we say so and score it anyway. Pulls, the scoring classes (scored,
not yet sensitive, INFO) and the verdict words are those of Section~\ref{sec:scoring}. Every measured number was checked on 2026-10-31
against its publication, on arXiv wherever possible. The exceptions are the CDF W mass, read from a fit paper that quotes it; the rubidium
value of $\alpha$, read from the journal abstract; Gasser and Leutwyler's 1982 value, read through a later thesis; and the PDG and CODATA values
that Section~\ref{sec:data} verified. The script's citation list gives, per source, what was read and the value.

\subsection{SM1: the spin ladder}\label{sec:sm-spin}

SM1 says that elementary quanta have spin 0, 1/2 or 1. The spin-2 wave belongs to the classical frame and has no quantum, and no elementary
spin-3/2 quantum exists. The test is the absence of what a richer ladder would add.
\begin{itemize}[nosep]
\item Supersymmetry, whose gravitino has spin 3/2, is bounded through its coloured partners. ATLAS excludes gluinos below 2.30\,TeV and
degenerate light squarks below 1.85\,TeV at 95\%~CL, in simplified models with a massless lightest neutralino~\cite{atlas2021a}
\chk{sm:SM1.1}.
\item A spin-2 Kaluza--Klein graviton of the Randall--Sundrum model is excluded below 4.5, 3.9 and 2.2\,TeV for $k/\bar M_{\rm Pl}=0.1$, 0.05
and 0.01~\cite{atlas2021b} \chk{sm:SM1.2}.
\item An excited top quark of spin 3/2 is excluded below 1.2\,TeV~\cite{cms2018} \chk{sm:SM1.3}.
\item The one scalar found has spin 0: every spin-2 alternative is excluded at more than 99.9\%~CL~\cite{atlas2015} \chk{sm:SM1.4}.
\end{itemize}
We found no gravitino-specific limit we could verify, and use none. Null searches bound masses; they cannot exclude a spin. The Standard
Model makes the same statement for its own content. \emph{Verdict: holds as a reading}, and it cannot become more.

\subsection{SM2: the Higgs vacuum grants every elementary rest mass}\label{sec:sm-higgs}

\paragraph{Couplings proportional to mass.} In the Standard Model a fermion of mass $m_f$ couples to the Higgs boson with $y_f=\sqrt2\,m_f/v$,
and a weak boson with $g_{hVV}=2m_V^2/v$, where $v=(\sqrt2G_F)^{-1/2}=246.22$\,GeV~\cite{mulan2013} \tagP. Hypothesis~1$'$ makes the grant the only
source of elementary rest mass, so it predicts exactly these couplings. The experiments measure coupling modifiers $\kappa_i$, which scale each
Standard Model coupling. The \emph{reduced couplings}
\begin{equation}\label{eq:sm-red}
y_f^{\rm red}=\kappa_f\frac{m_f}{v},\qquad y_V^{\rm red}=\sqrt{\kappa_V}\,\frac{m_V}{v}
\end{equation}
then lie on one line, $y^{\rm red}=m/v$, of slope 1 on logarithmic axes~\cite{atlas2022,cms2022a}. To test slope and intercept together we fit
\begin{equation}\label{eq:sm-fit}
\ln\kappa_{{\rm eff},i}=a+(s-1)\ln\frac{m_i}{m_*},\qquad \kappa_{\rm eff}=\kappa_f\ \text{or}\ \sqrt{\kappa_V},\qquad v_{\rm fit}=v\,e^{-a},
\end{equation}
by weighted least squares with the symmetrised errors of each $\kappa$. The pivot $m_*$ makes $a$ and $s$ uncorrelated. The prediction is $s=1$
and $v_{\rm fit}=v$.

The data are the CMS modifiers of the fit with effective photon, gluon and $Z\gamma$ couplings: $\kappa_W=1.02\pm0.08$, $\kappa_Z=1.04\pm0.07$,
$\kappa_t=1.01^{+0.11}_{-0.10}$, $\kappa_b=0.99^{+0.17}_{-0.16}$, $\kappa_\tau=0.92\pm0.08$ and $\kappa_\mu=1.12^{+0.21}_{-0.22}$~\cite{cms2022a}. ATLAS gives its
per-particle values only in figures~\cite{atlas2022}; we did not read them off and use its two-parameter fit instead. Fermion masses enter at
the electroweak scale, $m_f=y_f(M_Z)\,v/\sqrt2$, from the $\overline{\rm MS}$ Yukawa couplings~\cite{antusch2025}: 168.4, 2.84, 1.730 and 0.1019\,GeV for $t$,
$b$, $\tau$ and $\mu$. This is close to ATLAS's convention of quark masses at the Higgs mass. The $W$ and $Z$ enter with their pole masses
\chk{sm:SM2a}.

The fit gives $s=1.009\pm0.017$ and $v_{\rm fit}=244.0\pm5.7$\,GeV at the pivot $m_*=57$\,GeV, with $\chi^2=1.25$ for 4 degrees of freedom. The
pulls are $+0.54\sigma$ and $-0.39\sigma$ \chk{sm:SM2a.4}, \chk{sm:SM2a.5}. With pole and PDG masses the slope is again $1.009\pm0.017$: the mass
convention moves each point along the line, not off it. Jointly, the six tree-level modifiers give $\chi^2=1.70$ for 6 against $\kappa=1$, with
diagonal errors taken on the side towards 1 ($p=0.95$) \chk{sm:SM2a.1}. The covariance is not published in the text, so the collaborations' own
likelihoods decide. CMS finds $p=0.28$ for all nine modifiers~\cite{cms2022a}. ATLAS finds $p=0.56$ (with $\kappa_c=\kappa_t$), 0.65 ($\kappa_c$
free) and 0.61 (effective loop couplings)~\cite{atlas2022} \chk{sm:SM2a.2}. ATLAS's two-parameter fit, $\kappa_V=1.035\pm0.031$ and
$\kappa_F=0.95\pm0.05$ with correlation 0.39, gives $\chi^2=3.72$ for 2 ($p=0.16$; ATLAS quotes 0.14) \chk{sm:SM2a.3}. The largest single
departure is the loop-induced $\kappa_{Z\gamma}=1.65^{+0.34}_{-0.37}$, $+1.76\sigma$, which is not a mass coupling \chk{sm:SM2a.kZgamma}.

\begin{figure}[t]
\centering
\includegraphics[width=\linewidth]{sm_couplings.pdf}
\caption{(a) Reduced Higgs couplings, Eq.~\eqref{eq:sm-red}, against mass at the electroweak scale (CMS 2022~\cite{cms2022a}; fermion masses
from~\cite{antusch2025}). Black: the prediction $m/v$; dashed: the fit of Eq.~\eqref{eq:sm-fit}. Green bar: the CMS 95\% interval for charm;
triangle: the ATLAS 95\% bound $|\kappa_c|<4.2$~\cite{atlas2025}; arrow: the electron bound $|\kappa_e|<240$~\cite{cms2023} above the predicted
value (open circle). (b) The CMS modifiers (68\%; green: loop-induced) and ATLAS's $\kappa_V$, $\kappa_F$. What to look for: six couplings over
3.2 decades of mass lie on the line; charm and the electron are bounded, not measured.}
\label{fig:sm-couplings}
\end{figure}

What is not measured is untested.
\begin{itemize}[nosep]
\item Charm: CMS's 95\% interval $1.1<|\kappa_c|<5.5$ just excludes 1, while ATLAS's newer bound $|\kappa_c|<4.2$ includes it~\cite{cms2022a,atlas2025}
\chk{sm:SM2a.6}.
\item The electron: $B(H\to ee)<3.0\times10^{-4}$ bounds $|\kappa_e|<240$, against a predicted branching fraction of about
$5\times10^{-9}$~\cite{cms2023} \chk{sm:SM2a.7}.
\item The light quarks are unmeasured. The neutrinos' Dirac couplings, $y_\nu=\sqrt2m_\nu/v\approx2.9\times10^{-13}$ for the heaviest neutrino at
the oscillation minimum, are out of reach of any collider \chk{sm:SM2a.8}, \chk{sm:SM2a.9}.
\end{itemize}
\emph{Verdict: holds} for $W$, $Z$, $t$, $b$, $\tau$ and $\mu$; \emph{untested} for the rest. The statement is the Standard Model's, and the data
that support it support both.

\paragraph{Dirac neutrinos, and the assumption they need.} The grant alone does not make neutrinos Dirac particles. The Standard Model
admits one operator of dimension five, $(L\Phi)(L\Phi)/\Lambda$, built from the lepton doublet $L$ and the Higgs doublet $\Phi$. Once the
Higgs field takes its vacuum value, it gives a Majorana mass $m_\nu=f\,v^2/\Lambda$~\cite{weinberg1979,ma2009}. That mass comes from the same
vacuum, so Hypothesis~1$'$ allows it. What forbids it is a conserved lepton number. The combination $B-L$ (baryon minus lepton number) is the
one the Standard Model conserves exactly, and a Majorana mass changes it by two \tagP. Dirac neutrinos therefore need one assumption more
\chk{sm:SM2b.0}:
\begin{itemize}[nosep]
\item \emph{Added reading (not derived).} $B-L$ is one of the counts that every reader sees (Axiom~3), like electric charge. The count of
a closed part changes only when something carrying it crosses the part's boundary, by a grant or a loss that leaves a record; no act inside
the part changes it. Neutrinoless double-beta decay would create two units of lepton number in a closed nucleus, with nothing carrying them
in, and is excluded.
\end{itemize}
The reading forbids every Majorana mass, whatever its source. It still allows the proton decay of SM4, since $p\to e^+\pi^0$ conserves
$B-L$. We label it an added reading: the axioms do not say which conserved quantities are counts \chk{sm:SM2b.7}.

Given this reading, neutrinos have Dirac masses from Yukawa couplings to right-handed partners, and neutrinoless double-beta decay does not
occur at any rate. Through light-neutrino exchange its amplitude is proportional to the effective Majorana mass $m_{\beta\beta}=|\sum_iU_{ei}^2m_i|$,
which is zero \tagP. The searches agree.
\begin{itemize}[nosep]
\item KamLAND-Zen finds $T_{1/2}(^{136}\mathrm{Xe})>3.8\times10^{26}$\,yr at 90\%~CL, or $m_{\beta\beta}<28$ to 122\,meV depending on the nuclear
matrix element~\cite{kamlandzen2025} \chk{sm:SM2b.1}.
\item GERDA, the MAJORANA DEMONSTRATOR and LEGEND-200 together find $T_{1/2}(^{76}\mathrm{Ge})>1.9\times10^{26}$\,yr, or $m_{\beta\beta}<75$ to
200\,meV~\cite{legend2025} \chk{sm:SM2b.2}.
\end{itemize}
Were the neutrinos Majorana, the oscillation parameters of NuFIT~6.0~\cite{esteban2024} would put $m_{\beta\beta}$ between 18.2 and 48.2\,meV
for the inverted ordering. For the normal ordering it would lie between 1.5 and 3.7\,meV at zero lightest mass (Fig.~\ref{fig:sm-neutrino}a).
LEGEND-1000 aims at a $3\sigma$ discovery at $1.3\times10^{28}$\,yr, or 9 to 21\,meV~\cite{legend2021}, and nEXO at
$1.35\times10^{28}$\,yr~\cite{nexo2022} \chk{sm:SM2b.3}. A signal in any of them would refute SM2.2. A null result would not establish Dirac
neutrinos, since a normal-ordered Majorana spectrum can hide below every planned reach. \emph{Verdict: live, given $B-L$ conservation}; it holds
against every search so far.

\begin{figure}[t]
\centering
\includegraphics[width=\linewidth]{sm_neutrino.pdf}
\caption{(a) The effective Majorana mass that Majorana neutrinos would have, for either ordering (NuFIT~6.0~\cite{esteban2024}; bands over both
Majorana phases), against the KamLAND-Zen bound~\cite{kamlandzen2025} and the LEGEND-1000 reach~\cite{legend2021}. With $B-L$ conserved, GUID
predicts $m_{\beta\beta}=0$, below the plot. (b) 95\% upper bounds on $\sum m_\nu$~\cite{elbers2025,calabrese2025,planck18} against the oscillation
minima (dashed: normal, dotted: inverted ordering). (c) $N_{\rm eff}$ measured~\cite{planck18,calabrese2025} against decoupled right-handed partners
(GUID, 3.044) and partners thermalised before the electroweak split. What to look for: a signal in (a) would refute GUID's Dirac neutrinos; the
tension in (b) is shared with every neutrino model.}
\label{fig:sm-neutrino}
\end{figure}

Two side tests concern the right-handed partners and the mass scale. With only Yukawa interactions the partners never come into equilibrium,
and their contribution to the radiation density is negligible~\cite{luo2020}. Hence $N_{\rm eff}=3.0440$~\cite{bennett2021}. Planck with baryon
acoustic oscillations finds $2.99\pm0.17$~\cite{planck18}, and ACT DR6 with Planck, lensing and BAO finds $2.86\pm0.13$~\cite{calabrese2025}. The pulls
are $-0.32\sigma$ and $-1.42\sigma$ \chk{sm:SM2b.4}, \chk{sm:SM2b.5}. Had the partners been thermalised in the start and decoupled above every
Standard Model threshold, entropy conservation gives each a share $(10.75/106.75)^{4/3}$, so $N_{\rm eff}=3.185$. That lies $2.50\sigma$ above the
ACT value \chk{sm:SM2b.5b}. It is not GUID's statement, but it tells the start not to equilibrate them.

Oscillations bound the sum of the masses from below: $\sum m_\nu\ge0.0588$\,eV for the normal and $\ge0.0989$\,eV for the inverted ordering
(NuFIT quotes 0.058 and 0.098\,eV). Cosmology bounds it from above. DESI's baseline 95\% bound, $\sum m_\nu<0.0642$\,eV, lies 0.0054\,eV \emph{above}
the normal-ordering minimum. Its Feldman--Cousins limit, 0.053\,eV, breaches that minimum, and its analysis with an effective mass allowed to
go negative is in $3\sigma$ tension with it~\cite{elbers2025} \chk{sm:SM2b.6}. The inverted ordering is disfavoured. The tension is the same for
Dirac and Majorana neutrinos and is $\Lambda$CDM's: \emph{inherited}. The same bound makes the known neutrinos less than 0.57\% of the dark
matter \chk{data:D9.3}.

\paragraph{Massless bosons and the W mass.} That the photon and the gluons stay massless is a reading of the Standard Model, with no number
scored \chk{sm:SM2c.1}. At tree level $m_W=gv/2$ and $m_Z=m_W/\cos\theta_W$ come within $-0.18\%$ and $+0.34\%$ of the measured masses
\chk{data:D8.3}, \chk{data:D8.4}. Beyond tree level GUID adds no loop of its own, so its $m_W$ is the Standard Model's. The global electroweak fit
without $m_W$ gives $80.3499\pm0.0056$\,GeV~\cite{deblas2022}. The new measurements agree with it:
\begin{itemize}[nosep]
\item ATLAS, $80.3665\pm0.0159$\,GeV, $+0.98\sigma$~\cite{atlas2024b};
\item CMS, $80.3602\pm0.0099$\,GeV, $+0.91\sigma$~\cite{cms2026};
\item LHCb, $80.354\pm0.032$\,GeV, $+0.13\sigma$~\cite{lhcb2022} \chk{sm:SM2d.1-SM2d.3}.
\end{itemize}
The combination without CDF, which shares data with them, is $+1.34\sigma$~\cite{lhctev2024} \chk{sm:SM2d.5}. CDF's $80.4335\pm0.0094$\,GeV~\cite{cdf2022}
lies $7.6\sigma$ above (Fig.~\ref{fig:sm-precision}a) \chk{sm:SM2d.4}. We reproduce the fit authors' $6.5\sigma$ against the 2022 average that
included CDF. \emph{Verdict: holds}. CDF is an \emph{inherited} tension, now isolated by three experiments that agree with the Standard Model and
with each other.

\subsection{SM3: only the Maxwell term varies}\label{sec:sm-mu}

\paragraph{The ratio $R$.} Hypothesis~\ref{h:info} with the action~\eqref{eq:action} lets the scalar change one coefficient, that of
$F_{\mu\nu}F^{\mu\nu}$. The strong and weak gauge couplings, the Yukawa couplings and the Higgs potential do not move. Below the electroweak
scale this shifts $1/\alpha$ by a constant, the same at every scale at one loop, because the running of $1/\alpha$ does not depend on $\alpha$
there. Hence $\delta\ln\alpha(\mu)=[\alpha(\mu)/\alpha(0)]\,\delta\ln\alpha(0)$. The proton-to-electron mass ratio $\mu=m_p/M_e$ then changes only through
electromagnetic contributions, and
\begin{equation}\label{eq:sm-R}
R\equiv\frac{\dd\ln\mu}{\dd\ln\alpha}=\frac{\dd\ln m_p}{\dd\ln\alpha}-\frac{\dd\ln M_e}{\dd\ln\alpha}.
\end{equation}
The value depends on what is held fixed while $\alpha$ moves, the \emph{matching}. We compute two.

\emph{Matching at the electroweak scale.} The Yukawa couplings and $v$ are held fixed at $\mu_0=M_Z$, where Part~I places the birth of $\alpha$;
the three-flavour QCD scale $\Lambda_{\rm QCD}$ is held fixed too. The electron's $\overline{\rm MS}$ mass runs in QED as
$\dd\ln m/\dd\ln\mu=-3\alpha(\mu)/2\pi$ at one loop, and its pole mass is $M_e=m_e(m_e)(1+\alpha/\pi)$ \tagP. These are the QCD relations with the colour
factor $4/3$ replaced by 1. With $m_e(\mu_0)=y_ev/\sqrt2$ fixed,
\begin{equation}\label{eq:sm-me}
\frac{\dd\ln M_e}{\dd\ln\alpha}=\frac{3}{2\pi\alpha(0)}\int_{m_e}^{\mu_0}\alpha(\mu)^2\,\dd\ln\mu+\frac{\alpha}{\pi}=0.0436+0.0023=0.0460 .
\end{equation}
The running $\alpha(\mu)$ is the one-loop leptonic running, scaled to the three-loop $\Delta\alpha_{\rm lep}(M_Z^2)=314.977\times10^{-4}$~\cite{steinhauser1998},
plus the hadronic $\Delta\alpha^{(5)}_{\rm had}(M_Z^2)=(276.0\pm1.0)\times10^{-4}$~\cite{davier2020} spread logarithmically from 0.5\,GeV to $M_Z$. It
reproduces $1/\alpha(M_Z)=128.937$ on shell. With $\alpha$ frozen at $\alpha(0)$, the first term would be $(3\alpha/2\pi)\ln(M_Z/m_e)$, with
$\ln(M_Z/m_e)=12.09$ \chk{sm:SM3}. This is the electron's electromagnetic self-energy: 4.6\% of its mass.

The proton carries $\alpha$ in two pieces. One is its electromagnetic self-energy $E_{\rm em}$, at fixed quark masses at 2\,GeV. The other is the
QED running of those quark masses between $\mu_0$ and 2\,GeV, weighted by the sigma terms $f_q=\sigma_q/m_N$:
\begin{equation}\label{eq:sm-mp}
\frac{\dd\ln m_p}{\dd\ln\alpha}=\frac{E_{\rm em}}{m_p}+\sum_{q=u,d,s}f_qQ_q^2\,\frac{3}{2\pi\alpha(0)}\int_{2\,{\rm GeV}}^{\mu_0}\alpha(\mu)^2\,\dd\ln\mu
=6.7\times10^{-4}+2.2\times10^{-4}.
\end{equation}
We take $E_{\rm em}=0.63$\,MeV, Gasser and Leutwyler's Born-level value from the measured form factors~\cite{gasser1982}, read
through~\cite{stern2008}, with $\pm0.40$\,MeV. This error covers their estimate of $\pm0.2$\,MeV for intermediate states. It also covers the spread of
the proton-neutron difference between the Cottingham evaluation, $0.58\pm0.16$\,MeV, and lattice values near 1.0\,MeV~\cite{gasser2021}. The
individual self-energy is not an observable on its own, and its scheme dependence is part of this error. The sigma-term fractions are the
lattice values $f_u=0.0142$, $f_d=0.0242$ and $f_s=0.0577$~\cite{borsanyi2020}. Scaling $f_u$ and $f_d$ up to the dispersive
$\sigma_{\pi N}=59.1\pm3.5$\,MeV~\cite{hoferichter2015} changes $R$ by $8\times10^{-5}$ \chk{sm:SM3}.

Together, $R=8.9\times10^{-4}-0.0460=-0.0451\pm0.0023$ (Table~\ref{tab:sm-budget}). Over the range of choices the value lies between $-0.0489$ and
$-0.0444$ \chk{sm:SM3.1}. Its sign is fixed: a larger $\alpha$ gives the electron more self-energy, in proportion, than the proton.

\begin{table}[t]
\centering\small
\caption{The error budget of $R$ under the electroweak matching, Eqs.~\eqref{eq:sm-R}--\eqref{eq:sm-mp}; quadrature sum $\pm0.0023$.}\label{tab:sm-budget}
\begin{tabular}{@{}lr@{}}
\toprule
source & $\delta R$ \\
\midrule
matching scale $\mu_0$ between $M_W$ and $v$ (half range) & $2.2\times10^{-3}$ \\
$\Lambda^{(3)}_{\rm QCD}$ fixed against $\alpha_s(M_Z)$ fixed (heavy-quark thresholds) & $6.2\times10^{-4}$ \\
proton electromagnetic self-energy, $\pm0.40$\,MeV & $4.3\times10^{-4}$ \\
two-loop QED anomalous dimension, $O(\alpha/\pi)$ of the electron term & $1.1\times10^{-4}$ \\
sigma terms, lattice against dispersive & $7.5\times10^{-5}$ \\
hadronic running model, onset 0.3 to 1\,GeV (half range) & $6.2\times10^{-5}$ \\
\bottomrule
\end{tabular}
\end{table}

\emph{Matching at low energy.} Many analyses of varying constants hold the low-energy masses fixed instead. Then the electron term of
Eq.~\eqref{eq:sm-me} vanishes. With every low-energy mass fixed, $R=E_{\rm em}/m_p=(+6.7\pm4.3)\times10^{-4}$; with only the electron's fixed,
$R=(+8.9\pm7.6)\times10^{-4}$ \chk{sm:SM3.1b}.

\emph{Which matching Hypothesis~\ref{h:info} selects.} Holding the low-energy electron mass fixed while $\alpha$ moves requires the Yukawa
coupling to move against it, $\dd\ln y_e/\dd\ln\alpha=-0.046$. That is a coupling of the scalar to the electron's mass term. Hypothesis~\ref{h:info}
lets the scalar multiply the Maxwell term and nothing else, so the mass terms are fixed where the Maxwell term is defined. The Maxwell term
exists only below the electroweak split, so they are fixed at the electroweak scale; the grant of Hypothesis~1$'$ is fixed there too. We
therefore take the electroweak matching as GUID's. The exact scale within the electroweak range, $M_W$ to $v$, is part of
Table~\ref{tab:sm-budget}.

\paragraph{The preregistered window.} The preregistration expected $R$ ``of order $10^{-3}$ to $10^{-2}$''. Under the electroweak matching the
derived value lies a factor of 4.5 above that window, and no scale between $M_W$ and $v$ brings it inside. The estimate behind the window
counted the proton's electromagnetic mass, 0.09\% of it, and missed the electron's, 4.6\%. Under the low-energy matching $R$ is of order
$10^{-3}$, at the lower edge of the window and within one standard deviation of it. The preregistration did not fix the matching. We score the
window as written: \emph{refuted under the electroweak matching} that Hypothesis~\ref{h:info} selects, by derivation rather than by data, and
\emph{inside it under the low-energy matching} \chk{sm:SM3.1}, \chk{sm:SM3.1b}.

What the window was meant to say survives under both. Unification models tie $\mu$ to $\alpha$ through $\Lambda_{\rm QCD}$ and give $|R|$ of 20 to
50. In the class studied by Ferreira et al., $\Delta\mu/\mu=[0.8R_\Lambda-0.3(1+S)]\Delta\alpha/\alpha$, which is $-24.3$ for their typical
$R_\Lambda\sim30$ and $S\sim160$~\cite{ferreira2012}. Calmet and Fritzsch find the nucleon mass changing about 40 times faster than
$\alpha$~\cite{calmet2002}. GUID's $|R|$ is 540 times smaller than the first, and at least 400 times smaller than any $|R|\ge20$ \chk{sm:SM3}.

\paragraph{The data.} Each model is a line through the origin, $\Delta\mu/\mu=R\,\Delta\alpha/\alpha$; the data are bounds near zero on both axes
(Fig.~\ref{fig:sm-mualpha}). The rows are \emph{not yet sensitive}: no variation has been seen, so they cannot tell GUID's line from the
Standard Model's constant constants.
\begin{itemize}[nosep]
\item Methanol absorption at $z=0.886$ gives $|\Delta\mu/\mu|\le1.1\times10^{-7}$ at $2\sigma$, over a lookback time of 7.5\,Gyr~\cite{kanekar2015}; we
read it as a Gaussian $0\pm0.55\times10^{-7}$. ESPRESSO gives $\Delta\alpha/\alpha=1.3\pm1.3\pm0.4$\,ppm at $z=1.15$~\cite{espresso} \chk{data:D1.6}. On
GUID's line the two combine to $\chi^2=0.51$ for one degree of freedom ($p=0.48$) and to $\Delta\alpha/\alpha=0.58\pm0.91$\,ppm. Under GUID the
molecular bound limits $\alpha$ about as tightly as the quasar does \chk{sm:SM3.2}.
\item On the unification line $R=-24.3$ the same data are equally consistent ($p=0.34$), but they then force $|\Delta\alpha/\alpha|<4.4\times10^{-9}$
at 95\% \chk{sm:SM3.3}.
\item Molecular hydrogen at $z=2.0$ to 4.2 gives $|\Delta\mu/\mu|<5\times10^{-6}$ at $3\sigma$~\cite{ubachs2016}. That is no constraint on $\alpha$ under
GUID ($<1.1\times10^{-4}$), and $2.1\times10^{-7}$ under unification \chk{sm:SM3.5}.
\item Today's clocks give $\dot\alpha/\alpha=(1.8\pm2.5)\times10^{-19}\,\mathrm{yr}^{-1}$~\cite{filzinger2023} \chk{data:D1.3} and
$\dot\mu/\mu=(-8\pm36)\times10^{-18}\,\mathrm{yr}^{-1}$ from ytterbium against caesium~\cite{lange}. Both lines pass ($p=0.82$ and 0.92) \chk{sm:SM3.4}.
\end{itemize}
The test that would decide: a variation of $\alpha$ at the ppm level with $|\Delta\mu/\mu|$ below $10^{-7}$ would exclude unification and agree with
GUID. A $\Delta\mu/\mu$ much larger than $0.045\,|\Delta\alpha/\alpha|$ would refute Hypothesis~\ref{h:info}. \emph{Verdict: live.}

\begin{figure}[t]
\centering
\includegraphics[width=\linewidth]{sm_mu_alpha.pdf}
\caption{Variations of $\mu$ against variations of $\alpha$. (a) Near $z\approx1$: the ESPRESSO 95\% band in $\alpha$ ($z=1.15$)~\cite{espresso}
and the methanol $2\sigma$ band in $\mu$ ($z=0.89$)~\cite{kanekar2015}; the lines are GUID's $R=-0.045$ (electroweak matching) and
unification's $R=-24$ (dashed; dotted: $-20$ and $-50$)~\cite{ferreira2012}. The low-energy matching's line would be nearly horizontal.
(b) Present drifts from clocks~\cite{filzinger2023,lange}. What a model allows is its line inside both bands: a segment of ppm length for GUID,
a point within $4\times10^{-9}$ of zero for unification at $z\approx1$.}
\label{fig:sm-mualpha}
\end{figure}

\paragraph{Lepton magnetic moments.} The scalar couples to a lepton only through the lepton's electromagnetic self-energy. Its strength is
$\lambda\,m_\ell/\bar M_{\rm P}$ times a factor below 0.05 (0.046 for the electron, Eq.~\eqref{eq:sm-me}); to photons it is $\lambda/\bar M_{\rm P}$. Any
contribution to $a_\ell=(g-2)/2$ needs two such insertions and one helicity flip. Hence $\Delta a_\ell=C\lambda^2m_\ell^2/\bar M_{\rm P}^2$, where $C$
collects loop factors and logarithms; we take $C\le10$, generously. At $|\lambda|<3.1\times10^{-5}$~\eqref{eq:lambdabound} this gives
$\Delta a_e\lesssim4\times10^{-52}$ and $\Delta a_\mu\lesssim2\times10^{-47}$. These lie 38 and 36 orders of magnitude below the experimental errors
\chk{sm:SM3.6}: \emph{derived}, and \emph{holds}, in a row that is not yet sensitive.

Nor can the scalar make $\alpha$ differ between laboratories. A body of electrostatic share $f$ sources $\phi=2\lambda f\,\Phi/c^2$, the coupling
behind Eq.~\eqref{eq:fifth}. So $\delta\alpha/\alpha=2\lambda^2f\,\delta\Phi/c^2\le10^{-20}$ for $f\le5\times10^{-3}$ and $\delta\Phi/c^2\le10^{-9}$, which
covers the Earth's surface potential and the annual change of the Sun's \chk{sm:SM3.7}. The two best determinations of $\alpha$ differ by
$1.2\times10^{-9}$, or $5.5\sigma$: $\alpha^{-1}=137.035999046(27)$ from caesium~\cite{parker2018} and $137.035999206(11)$ from
rubidium~\cite{morel2020}. That difference is $10^{11}$ times what the scalar allows. It is experimental, or lies in the inputs of the
Standard Model; it is not GUID's.

The electron's moment, $g/2=1.00115965218059(13)$, gives with the Standard Model $\alpha^{-1}=137.035999166(15)$~\cite{fan}. Against rubidium that
is $+2.15\sigma$ ($a_e^{\rm exp}-a_e^{\rm SM}=+0.34\times10^{-12}$), and against caesium $-3.89\sigma$ ($-1.02\times10^{-12}$). Both FAIL at 95\%
\chk{sm:SM3.8}, \chk{sm:SM3.9}. They are one \emph{inherited} tension, in the determinations of $\alpha$, seen twice. The muon is clean: the
world average after Fermilab's final result, $a_\mu=116\,592\,071.5(14.5)\times10^{-11}$~\cite{muongtwo2025}, lies $+0.60\sigma$ from the 2025
Standard Model value $116\,592\,033(62)\times10^{-11}$~\cite{aliberti2025} \chk{sm:SM3.10} (Fig.~\ref{fig:sm-precision}b,c). \emph{Verdict} of the
side prediction: \emph{derived}; \emph{holds}.

\begin{figure}[t]
\centering
\includegraphics[width=\linewidth]{sm_precision.pdf}
\caption{(a) The W mass: measurements~\cite{atlas2024b,cms2026,lhcb2022,lhctev2024,cdf2022} against the Standard Model fit without
$m_W$~\cite{deblas2022} (band, $1\sigma$). (b) $\alpha^{-1}$ from caesium~\cite{parker2018}, rubidium~\cite{morel2020}, the electron's
moment~\cite{fan} and CODATA 2022~\cite{codata2022}. (c) The muon anomaly, world average~\cite{muongtwo2025} and Standard Model~\cite{aliberti2025}. What
to look for: CDF alone departs in (a); in (b) the electron's moment sits between two determinations of $\alpha$ that disagree with each other;
(c) agrees.}
\label{fig:sm-precision}
\end{figure}

\subsection{SM4: no unification within causality}\label{sec:sm-unif}

Hypothesis~6$'$ puts any unification of the strong with the electroweak force before the start. Within causality, then, the Standard Model
couplings run with Standard Model content alone should not meet at one point. We run the three couplings at two loops, with
$\alpha_1=\frac53\alpha_Y$ normalised for grand unification:
\begin{equation}\label{eq:sm-rge}
\frac{\dd\alpha_i^{-1}}{\dd\ln\mu}=-\frac{b_i}{2\pi}-\sum_j\frac{b_{ij}\alpha_j}{8\pi^2}+\frac{d_i\,y_t^2}{32\pi^3},
\end{equation}
with $b=(41/10,-19/6,-7)$, $b_{ij}=\bigl(\begin{smallmatrix}199/50&27/10&44/5\\9/10&35/6&12\\11/10&9/2&-26\end{smallmatrix}\bigr)$, $d=(17/10,3/2,2)$ and
$16\pi^2\,\dd y_t/\dd\ln\mu=y_t(\tfrac92y_t^2-8g_3^2-\tfrac94g_2^2-\tfrac{17}{20}g_1^2)$~\cite{buttazzo2013}. For contrast, the minimal supersymmetric
extension takes over above a common threshold $M_{\rm SUSY}=2$\,TeV. It has $b=(33/5,1,-3)$,
$b_{ij}=\bigl(\begin{smallmatrix}199/25&27/5&88/5\\9/5&25&24\\11/5&9&14\end{smallmatrix}\bigr)$ and $d=(26/5,6,4)$~\cite{martin1994}, its own top-Yukawa
equation~\cite{martin1997}, $y_t\to y_t/\sin\beta$ and $\tan\beta=10$. The inputs are $1/\hat\alpha(M_Z)=127.950\pm0.017$~\cite{erler2017},
$\sin^2\hat\theta_W=0.23129\pm0.00004$~\cite{navas2024}, $\alpha_s(M_Z)=0.1180\pm0.0009$~\cite{antusch2025,denterria2024} and $y_t(M_Z)=0.967$~\cite{antusch2025},
so $\alpha_i^{-1}(M_Z)=59.01$, 29.59 and 8.47. As the measure of mismatch we take the smallest spread $\Delta=\max_i\alpha_i^{-1}-\min_i\alpha_i^{-1}$
over $10^{12}$ to $10^{19}$\,GeV, which lies at a pairwise crossing. Its error combines 200 draws of the inputs (fixed seed) with a theory allowance.
The allowance is the whole two-loop correction, $|\Delta_{\rm 2\,loop}-\Delta_{\rm 1\,loop}|=0.47$, a deliberately generous stand-in for missing higher
orders and thresholds \chk{sm:SM4}.

In the Standard Model $\alpha_1$ meets $\alpha_2$ at $1.08\times10^{13}$\,GeV, $\alpha_1$ meets $\alpha_3$ at $1.71\times10^{14}$\,GeV, and $\alpha_2$
meets $\alpha_3$ at $2.93\times10^{16}$\,GeV (Fig.~\ref{fig:sm-running}). The smallest spread is $\Delta=3.20$ at $1.7\times10^{14}$\,GeV, against 0.045
from the inputs and 0.47 from the allowance: $6.8\sigma$ from a single point \chk{sm:SM4.1}. Put differently, a single meeting point would need
$\alpha_s(M_Z)=0.0738$, 37\% below the measured value and $49\sigma$ from it \chk{sm:SM4.2}. The supersymmetric contrast meets at
$1.15\times10^{16}$\,GeV with $1/\alpha_G=25.66$, and needs $\alpha_s(M_Z)=0.1181$, within the input error, before any threshold correction
\chk{sm:SM4.3}. A unification within causality would need such content beyond the Standard Model; superpartners near 2\,TeV are bounded by
the searches of SM1, not excluded.

Proton decay is the remaining trace of a unification before the start. In Part~I's estimate $\tau\sim M_X^4/(\alpha_G^2m_p^5)$, Super-Kamiokande's
$\tau/B(p\to e^+\pi^0)>2.4\times10^{34}$\,yr~\cite{superk} requires $M_X>4.8\times10^{15}$\,GeV for $\alpha_G=1/40$, and $4.7\times10^{15}$\,GeV for the
Standard Model's $\alpha_G=1/42.2$ \chk{sm:SM4.4}, \chk{data:D8.1}. $M_X$ is free, so the bound tests nothing GUID fixes. A unification at the Standard
Model's $\alpha_1=\alpha_2$ crossing would give $\tau\sim7\times10^{23}$\,yr, $3.5\times10^{10}$ times too short \chk{sm:SM4.5}. \emph{Verdict: holds as
a reading}: that the couplings do not meet is the Standard Model's own statement. Hypothesis~6$'$ itself, a unification before the start, is
\emph{untested}; Hyper-Kamiokande~\cite{hyperk} is its test.

\begin{figure}[t]
\centering
\includegraphics[width=\linewidth]{sm_running.pdf}
\caption{(a) The inverse couplings at two loops, Eq.~\eqref{eq:sm-rge}, in the Standard Model (solid) and with supersymmetric content above
2\,TeV (dashed), from the measured values at $M_Z$~\cite{erler2017,navas2024,antusch2025}. (b) Their spread $\max_i\alpha_i^{-1}-\min_i\alpha_i^{-1}$;
dots: the smallest spread. What to look for: the Standard Model's three crossings span three decades in energy and never close to a point.}
\label{fig:sm-running}
\end{figure}

\subsection{SM5: the coupling at the split}\label{sec:sm-split}

Part~I says that $\alpha$ is born near $1/128$ at the Z mass and runs down to $1/137$. From $\alpha(0)$~\cite{codata2022} and the measured vacuum
polarisation, $1/\alpha(M_Z)=\alpha(0)^{-1}[1-\Delta\alpha_{\rm lep}-\Delta\alpha_{\rm had}^{(5)}]=128.937\pm0.014$ on shell~\cite{steinhauser1998,davier2020},
$+0.73\%$ from 128 \chk{sm:SM5.1}. In the $\overline{\rm MS}$ scheme (modified minimal subtraction) the value is $127.950\pm0.017$~\cite{erler2017},
$-0.04\%$ from 128 \chk{data:D8.5}. The scheme offset implied by Erler's paired values~\cite{erler1999}, 0.98, turns our on-shell value into 127.96
\chk{sm:SM5.2}. The two schemes differ by about one unit in $1/\alpha$, which is as precisely as ``near $1/128$'' can be meant. Part~I took the value from this physics, so the row is a consistency row,
not counted as a success (Section~\ref{sec:scoring}). The split itself
is the lattice crossover at $T_c=159.5\pm1.5$\,GeV, with the measured Higgs mass above the end point of the first-order line \chk{data:D8.6}.
\emph{Verdict: holds}, in either scheme, as the Standard Model's.

\subsection{Choices, caveats and the count}\label{sec:sm-end}

\paragraph{Choices made after seeing the data.}
\begin{itemize}[nosep]
\item ATLAS publishes its per-particle coupling modifiers only in figures, so the mass-coupling fit uses the CMS values; ATLAS enters through its
two-parameter fit and its $p$-values.
\item The two rows of the electron's moment are scored separately, although they show one tension. This was decided after both were seen to fail.
\item The smallest spread of the couplings is evaluated at their crossings rather than on a grid, after the grid minimum was seen to sit at a
kink (a change below 0.02).
\item The argument that Hypothesis~\ref{h:info} selects the electroweak matching was written after $R$ had been computed, and after a reader
pointed out that the low-energy matching puts $R$ inside the window. The electroweak baseline itself was chosen before computing $R$.
\item The Dirac step is made conditional on $B-L$ conservation after a reader pointed out the dimension-five operator. The data and their
outcome are unchanged.
\end{itemize}
Chosen before the comparison: the error of the proton's self-energy, the supersymmetric threshold of 2\,TeV with $\tan\beta=10$, and reading the
methanol $2\sigma$ bound as a Gaussian centred on zero. No verdict depends on these choices, except the window of SM3, which depends on the matching
as stated above.

\paragraph{Caveats.}
\begin{itemize}[nosep]
\item The coupling fit ignores correlations between the $\kappa$, which are not published in the text. The collaborations' own $p$-values, which
include them, agree with the verdict.
\item The proton's electromagnetic self-energy is not an observable on its own; its scheme dependence is covered only by the error we assign.
\item The matching of the electron's mass at the electroweak scale has electroweak corrections of order $\alpha/\pi$ that we bound but do not
compute.
\item The methanol and quasar data lie at different redshifts and along different lines of sight; combining them assumes a slowly varying
scalar. The clock values of $\alpha$ and $\mu$ share data, and the caesium-based $\mu$ assumes a model for the caesium nuclear $g$-factor.
\item The two-loop running omits threshold corrections at the top mass, the conversion to the supersymmetric scheme in the contrast model, and all
corrections at a unification scale. The theory allowance is a deliberately large stand-in for them.
\item The Standard Model value of $m_W$ was fitted with the 2022 top mass. CDF's $m_W$ was read from a paper quoting it, and the rubidium $\alpha$
from the journal abstract.
\end{itemize}

\begin{table}[t]
\centering\small
\caption{The preregistered statements and their verdicts. Rows are those of \texttt{standard\_model.json}; ``data:'' rows are those of
Section~\ref{sec:data}.}\label{tab:sm-verdicts}
\begin{tabularx}{\linewidth}{@{}lLLl@{}}
\toprule
& statement & verdict & rows \\
\midrule
SM1 & no elementary spin-3/2 or spin-2 quantum & holds as a reading (the Standard Model says the same) & SM1.1--1.4 \\
SM2.1 & Higgs couplings $\propto$ mass, all $\kappa=1$ & holds for $W,Z,t,b,\tau,\mu$; untested for $c$, $e$, $u$, $d$, $s$, $\nu$ & SM2a.1--9 \\
SM2.2 & Dirac neutrinos, $m_{\beta\beta}=0$ & live, given $B-L$ conservation (an added reading); holds against every search & SM2b.0--7 \\
SM2.3 & photon and gluons massless & holds as a reading & SM2c.1 \\
SM2.4 & $m_W=gv/2$, $\rho=1$, with Standard Model loops & holds; CDF 2022 inherited & SM2d.1--6 \\
SM3 & only $\alpha$ varies; $R$ ``of order $10^{-3}$ to $10^{-2}$'' & window refuted under the electroweak matching ($R=-0.045\pm0.002$), inside it under the
low-energy matching; live against data & SM3.1--5 \\
SM3 side & scalar's share of lepton $g-2$ negligible & derived; holds; $a_e$ limited by the inherited $\alpha$ discrepancy & SM3.6--10 \\
SM4 & no unification within causality & holds as a reading (the Standard Model's statement); Hypothesis~6$'$ untested & SM4.1--5 \\
SM5 & $\alpha$ born near $1/128$ at $m_Z$, runs to $1/137$ & holds & SM5.1--4 \\
\bottomrule
\end{tabularx}
\end{table}

\paragraph{The count.} Of 29 scored rows, 25 agree at 95\% \chk{sm:summary}. Twenty-two test statements GUID shares with the Standard Model or
$\Lambda$CDM: 19 agree, and the 3 FAILs are inherited (CDF's W mass and the electron's moment against either value of $\alpha$). Three are
GUID-specific and sensitive: the two neutrinoless double-beta searches agree, and the window of $R$ fails under the electroweak matching. Four
GUID-specific rows are not yet sensitive and agree (SM3.2, SM3.4 to SM3.6). Table~\ref{tab:sm-verdicts} lists the verdicts. All numbers in this
section come from \texttt{docs/simulation/standard\_model.py} (32 of 32 checks, about five seconds on one core) or, where marked ``data:'', from
the script of Section~\ref{sec:data}.

\paragraph{Verdict.} GUID's statements about the Standard Model hold wherever the data reach them, and almost everywhere they hold as the Standard
Model's own: couplings proportional to mass over three decades, the W mass at loop level, couplings that do not meet without new content, and
$\alpha$ near $1/128$ at the split. The three FAILs among them are inherited. Dirac neutrinos need one assumption beyond the grant, that $B-L$ is a
conserved count; given it, no neutrinoless double-beta decay occurs at any rate, a live prediction. The variation of $\mu$ that
Hypothesis~\ref{h:info} allows is $R=-0.045\pm0.002$ times that of $\alpha$ under the matching the hypothesis selects, outside the preregistered
window; under the low-energy matching it would be inside. Either way it is at least 400 times smaller than unification requires.
